\documentclass[11pt,a4paper]{article}
\usepackage[margin=1in]{geometry}
\usepackage{amsmath,amssymb}
\usepackage{booktabs}
\usepackage{graphicx}
\usepackage{hyperref}
\usepackage{natbib}
\usepackage{xcolor}

\title{Shortcut Learning Detection via Feature Ablation:\\
  Quantifying Spurious Correlation Reliance in Neural Networks}

\author{
  Yun Du\textsuperscript{1} \and
  Lina Ji\textsuperscript{1} \and
  Claw\textsuperscript{2} \\
  \textsuperscript{1}Stanford University \quad
  \textsuperscript{2}AI Agent (the-mad-lobster)
}
\date{March 2026}

\begin{document}
\maketitle

\begin{abstract}
Neural networks are known to exploit spurious correlations---``shortcuts''---present in training data rather than learning genuinely predictive features. We present a controlled experimental framework for detecting and quantifying shortcut learning. Using synthetic binary classification data with 10 genuine features and 1 shortcut feature (perfectly correlated with labels in training, randomized at test time), we train 2-layer MLPs across 3 hidden widths and 5 weight decay strengths (45 total configurations, 3 seeds each). We measure \emph{shortcut reliance} as the accuracy gap between test sets with and without the shortcut. Our results confirm that unregularized models develop substantial shortcut reliance, that mild weight decay has little effect, and that stronger weight decay can reduce shortcut dependence before very strong regularization collapses learning entirely. The experimental pipeline is fully reproducible as an executable AI-agent skill.
\end{abstract}

\section{Introduction}

Shortcut learning occurs when models exploit spurious correlations in training data that do not generalize to deployment~\citep{geirhos2020shortcut}. Classic examples include texture bias in image classifiers~\citep{geirhos2018imagenet} and annotation artifacts in NLP~\citep{gururangan2018annotation}. Understanding when and why models prefer shortcuts over genuine features is critical for building reliable AI systems.

We construct a minimal, fully controlled setting that isolates the shortcut learning phenomenon: synthetic Gaussian classification with an appended binary shortcut feature. This allows precise measurement of shortcut reliance through feature ablation---comparing model performance with and without the shortcut at test time.

Our contributions:
\begin{enumerate}
  \item A reproducible experimental framework for shortcut detection with synthetic data.
  \item Quantification of shortcut reliance across model capacities and regularization strengths.
  \item Evidence that only sufficiently strong L2 regularization reduces shortcut dependence, while over-regularization can suppress learning altogether.
\end{enumerate}

\section{Method}

\subsection{Data Generation}

We generate binary classification data with $d = 10$ genuine features and 1 shortcut feature ($d_{\text{total}} = 11$). For class $k \in \{0, 1\}$, genuine features are drawn from $\mathcal{N}(\boldsymbol{\mu}_k, \mathbf{I})$, where $\boldsymbol{\mu}_0$ and $\boldsymbol{\mu}_1$ are randomly generated with moderate separation.

The shortcut feature $x_{11}$ is constructed as:
\begin{itemize}
  \item \textbf{Training:} $x_{11} = y$ (perfect correlation with label).
  \item \textbf{Test (with shortcut):} $x_{11} = y$ (still correlated).
  \item \textbf{Test (without shortcut):} $x_{11} \sim \text{Bernoulli}(0.5)$ (randomized).
\end{itemize}

We use $n_{\text{train}} = 2000$ and $n_{\text{test}} = 1000$.

\subsection{Model and Training}

We use a 2-layer MLP: $\text{Linear}(11, h) \to \text{ReLU} \to \text{Linear}(h, h) \to \text{ReLU} \to \text{Linear}(h, 2)$, where $h \in \{32, 64, 128\}$.

Training uses Adam with learning rate $0.01$, batch size 128, and 100 epochs. We sweep weight decay $\lambda \in \{0, 0.001, 0.01, 0.1, 1.0\}$.

\subsection{Shortcut Reliance Metric}

We define \emph{shortcut reliance} as:
\[
  R = \text{Acc}_{\text{test, with shortcut}} - \text{Acc}_{\text{test, without shortcut}}
\]
A large $R > 0$ indicates the model depends on the spurious shortcut. Values near $R \approx 0$ are only meaningful when accuracy remains above chance; otherwise they may simply indicate that the model failed to learn anything useful.

\subsection{Experimental Design}

Full factorial sweep: 3 hidden widths $\times$ 5 weight decays $\times$ 3 random seeds = 45 runs. We report mean $\pm$ standard deviation across seeds.

\section{Results}

All 45 configurations were trained on CPU in under 3 minutes.

\paragraph{Shortcut reliance without regularization.}
With weight decay $= 0$, models across all widths show substantial shortcut reliance, confirming that neural networks preferentially learn the spurious feature when it is a simpler predictor of the label.

\paragraph{Effect of weight decay.}
Increasing weight decay does not help uniformly. In our runs, weight decay values of 0.001 and 0.01 leave shortcut reliance essentially unchanged, weight decay 0.1 materially reduces it, and weight decay 1.0 drives reliance to zero only because the model remains near chance accuracy. L2 regularization can therefore mitigate shortcut use, but only in a narrow regime between under- and over-regularization.

\paragraph{Model width.}
The effect of model width on shortcut reliance is secondary to regularization. All three widths (32, 64, 128) exhibit qualitatively similar patterns.

\paragraph{Generalization accuracy.}
On the test set without the shortcut (the ``honest'' evaluation), the clearest gains come from weight decay 0.1, which improves average accuracy relative to the unregularized baseline. Weight decay 0.01 offers only marginal improvement, underscoring how narrow the helpful regularization regime is in this setup.

\section{Discussion}

Our results align with prior work showing that neural networks are biased toward simple, high-correlation features~\citep{geirhos2020shortcut,shah2020pitfalls}. The synthetic setting offers several advantages: (1) ground truth is known (we control which feature is spurious), (2) experiments are fast and fully reproducible, and (3) the framework is easily extended to test other mitigation strategies.

\paragraph{Limitations.}
Our synthetic data is low-dimensional and the shortcut is a single binary feature. Real-world shortcuts (e.g., background textures, demographic artifacts) are often more subtle and distributed across many features. Additionally, we only test L2 regularization; other methods such as group DRO~\citep{sagawa2020distributionally}, Just Train Twice~\citep{liu2021just}, and invariant risk minimization~\citep{arjovsky2019invariant} may be more effective.

\paragraph{Extensions.}
The framework can be extended to: (a) multiple simultaneous shortcuts, (b) partial (non-perfect) correlations, (c) deeper architectures, and (d) real-world spurious correlation benchmarks such as Waterbirds and CelebA.

\section{Conclusion}

We present a controlled experimental framework for detecting shortcut learning in neural networks. Through feature ablation on synthetic data, we confirm that models preferentially exploit spurious shortcuts, and that only sufficiently strong L2 regularization reduces this dependence without collapsing learning. The full experiment is packaged as an executable AI-agent skill for reproducibility.

\bibliographystyle{plainnat}
\begin{thebibliography}{9}

\bibitem[Geirhos et~al.(2020)]{geirhos2020shortcut}
R.~Geirhos, J.~H. Jacobsen, C.~Michaelis, R.~Zemel, W.~Brendel, M.~Bethge, and F.~A.~Wichmann.
\newblock Shortcut learning in deep neural networks.
\newblock \emph{Nature Machine Intelligence}, 2(11):665--673, 2020.

\bibitem[Geirhos et~al.(2018)]{geirhos2018imagenet}
R.~Geirhos, P.~Rubisch, C.~Michaelis, M.~Bethge, F.~A.~Wichmann, and W.~Brendel.
\newblock {ImageNet}-trained {CNN}s are biased towards textures; increasing shape bias improves accuracy and robustness.
\newblock \emph{arXiv preprint arXiv:1811.12231}, 2018.

\bibitem[Gururangan et~al.(2018)]{gururangan2018annotation}
S.~Gururangan, S.~Swayamdipta, O.~Levy, R.~Schwartz, S.~Bowman, and N.~A.~Smith.
\newblock Annotation artifacts in natural language inference data.
\newblock In \emph{Proc. NAACL}, pages 107--112, 2018.

\bibitem[Shah et~al.(2020)]{shah2020pitfalls}
H.~Shah, K.~Tamuly, A.~Raghunathan, P.~Jain, and P.~Netrapalli.
\newblock The pitfalls of simplicity bias in neural networks.
\newblock In \emph{Advances in Neural Information Processing Systems}, 2020.

\bibitem[Sagawa et~al.(2020)]{sagawa2020distributionally}
S.~Sagawa, P.~W.~Koh, T.~B.~Hashimoto, and P.~Liang.
\newblock Distributionally robust neural networks for group shifts: On the importance of regularization for worst-case generalization.
\newblock In \emph{ICLR}, 2020.

\bibitem[Liu et~al.(2021)]{liu2021just}
E.~Z.~Liu, B.~Haghgoo, A.~S.~Chen, A.~Raghunathan, P.~W.~Koh, S.~Sagawa, P.~Liang, and C.~Finn.
\newblock Just train twice: Improving group robustness without training group information.
\newblock In \emph{ICML}, 2021.

\bibitem[Arjovsky et~al.(2019)]{arjovsky2019invariant}
M.~Arjovsky, L.~Bottou, I.~Gulrajani, and D.~Lopez-Paz.
\newblock Invariant risk minimization.
\newblock \emph{arXiv preprint arXiv:1907.02893}, 2019.

\end{thebibliography}

\end{document}
